\documentclass[../../main.tex]{subfiles}
\begin{document}

\begin{center}
  {\fontsize{48pt}{58pt}\selectfont \textbf{1978}\normalsize\textbf{年度}}
\end{center}

\begin{table}[htb]
  \centering
  \caption{出題テーマと難易度}
  \begin{tabular}{|c|p{9cm}|c|}
    \hline
    \textbf{問題番号} & \textbf{テーマ} & \textbf{難易度} \\
    \hline
    第1問 & 恒等式と係数比較 & \\
    \hline
    第2問 & 整数問題 & \\
    \hline
    第3問 & ベクトルと直線 & \\
    \hline
    第4問 & 数列と漸化式 & \\
    \hline
    第5問 & 積分と最小値 & \\
    \hline
    第6問 & 積分と最大値 & \\
    \hline
  \end{tabular}
\end{table}

\begin{table}[htb]
  \centering
  \caption{解答}
  \renewcommand{\arraystretch}{1.4}
  \begin{tabular}{|c|c|p{9cm}|}
    \hline
    \textbf{問題番号} & \textbf{小問} & \textbf{答え} \\
    \hline
    第1問 & - & $D=\dfrac{1}{16}$ \\
    \hline
    第2問 & (1) & 証明問題 \\
    \cline{2-3}
    & (2) & $(a,b,c)=(2,3,5)$ \\
    \hline
    第3問 & (1) & 証明問題 \\
    \cline{2-3}
    & (2) & $\dfrac{1}{2}\le t\le1$ \\
    \hline
    第4問 & - & $S=6p^3+\dfrac{9}{2}p^2-\dfrac{13}{2}p$ \\
    \hline
    第5問 & - & $(a,b,c,d)=\left(0,\dfrac{8}{7},-1,-\dfrac{1}{7}\right)$ \\
    \hline
    第6問 & - & 最大値 $e+\dfrac{1}{e}$ \\
    \hline
  \end{tabular}
\end{table}

\paragraph{試験形態}
6問．試験時間は要確認．

\end{document}
